% ============================================================
%  §2.3 Μηχανισμοί φθοράς -- κείμενο
% ============================================================

Η φθορά αγωνιστικού ελαστικού δεν αντιστοιχεί σε ένα μόνο φυσικό φαινόμενο. Στην πράξη
συνυπάρχουν τρεις ξεχωριστοί μηχανισμοί αποδόμησης του πέλματος, ο καθένας από τους
οποίους κυριαρχεί σε διαφορετικό θερμοκρασιακό εύρος και σε διαφορετικές συνθήκες
ολίσθησης. Η κατανόηση της φύσης τους είναι απαραίτητη τόσο για τη σωστή μοντελοποίηση όσο
και για την ερμηνεία των αποτελεσμάτων.

% ---
\subsubsection{Τριβολογική φθορά (Abrasion)}
\label{subsubsec:abrasion}
% ---

Η τριβολογική φθορά -- γνωστή και ως \emph{abrasion} -- αποτελεί τον «κανονικό» ρυθμό
αποδόμησης που εμφανίζεται όταν το ελαστικό λειτουργεί εντός του θερμικού παραθύρου. Ο
μηχανισμός είναι κλασικά τριβολογικός: οι ανωμαλίες της ασφάλτου (micro-asperities)
«ξύνουν» την επιφάνεια του πέλματος σε κάθε επαφή, αφαιρώντας μικροσκοπικά τμήματα
ελαστομερούς. Σύμφωνα με το νόμο Archard \cite{Archard1953}, ο ρυθμός φθοράς είναι
ανάλογος της ενέργειας τριβής στο τύπωμα, η οποία αντιστοιχεί άμεσα στον όρο $Q_1$:

\begin{equation}
    W_p = W_{p1} \cdot \left(\frac{Q_1}{Q_{\text{ref}}}\right)^{W_{p2}}
    \label{eq:Wp}
\end{equation}

όπου $Q_{\text{ref}}$ μια τιμή ισχύος αναφοράς και $W_{p1}$, $W_{p2}$ εμπειρικές
παράμετροι. Η εκθετική μορφή αντανακλά την επιτάχυνση της φθοράς σε συνθήκες υψηλής
ολίσθησης: το πρόβλημα δεν είναι γραμμικό -- διπλάσιο slip δεν σημαίνει διπλάσια φθορά αλλά
πολύ μεγαλύτερη αν $W_{p2} > 1$.

% ---
\subsubsection{Κόκκοι (Graining)}
\label{subsubsec:graining}
% ---

Το \emph{graining} είναι ένα φαινόμενο που εμφανίζεται όταν το ελαστικό λειτουργεί σε
θερμοκρασία \emph{κάτω} από το κατώτατο όριο του thermal window. Σε αυτές τις συνθήκες, το
ελαστόμερο είναι σχετικά σκληρό και δεν επιτυγχάνεται επαρκής χημική προσκόλληση με το
οδόστρωμα. Αντί να τριβόλεωνεται ομαλά, η επιφάνεια τείνει να αποκολλά μικρά τεμάχια
(«κόκκοι»), που στη συνέχεια επικολλούνται εκ νέου στο πέλμα αλλά με τυχαίο τρόπο -- η
επιφάνεια αποκτά χαρακτηριστική τραχιά, κοκκώδη υφή που μειώνει περαιτέρω το contact
patch:

\begin{equation}
    W_g = W_{g1} \cdot (T_{\text{tp}} - T)^{W_{g2}}, \quad T < T_{\text{tp}}
    \label{eq:Wg}
\end{equation}

Το graining είναι εν μέρει \emph{αναστρέψιμο}: αν το ελαστικό φτάσει στο thermal window, η
συσσωρευμένη θερμότητα μπορεί να «ανακυκλώσει» τους κόκκους και να αποκαταστήσει μέρος της
επιφάνειας. Στην πράξη παρατηρείται ότι οδηγοί που καταφέρνουν να ζεστάνουν τα ελαστικά
γρήγορα στην out-lap περιορίζουν σημαντικά το graining.

% TODO: Αν έχεις φωτογραφίες/εικόνες ελαστικών με graining vs. blistering,
% η παρουσίασή τους εδώ θα ήταν εξαιρετικά εποπτική για τον αναγνώστη.

% ---
\subsubsection{Φουσκάλες (Blistering)}
\label{subsubsec:blistering}
% ---

Το \emph{blistering} είναι ο αντίποδας του graining: εμφανίζεται όταν το ελαστικό
υπερθερμαίνεται, δηλαδή λειτουργεί \emph{πάνω} από το ανώτατο όριο του thermal window. Η
εξαιρετικά υψηλή θερμοκρασία τήκει τοπικά τα στρώματα ελαστομερούς, δημιουργώντας
εγκλωβισμένους κόκκους αερίων που σχηματίζουν φυσαλίδες (blisters) κάτω από την επιφάνεια.
Όταν οι φυσαλίδες αυτές σπάσουν, εκτινάσσουν κομμάτια πέλματος, καταστρέφοντας
ανεπανόρθωτα την επιφάνεια:

\begin{equation}
    W_b = W_{b1} \cdot (T - T_{\text{tp}})^{W_{b2}}, \quad T > T_{\text{tp}}
    \label{eq:Wb}
\end{equation}

Σε αντίθεση με το graining, το blistering είναι \emph{μη αναστρέψιμο}: μόλις δημιουργηθούν
φυσαλίδες, το ελαστικό έχει ήδη υποστεί δομική βλάβη. Αυτός είναι ο λόγος που οι μηχανικοί
αγώνα παρακολουθούν στενά την εξέλιξη της θερμοκρασίας ελαστικών κατά τη διάρκεια του
αγώνα και προσαρμόζουν τη στρατηγική (short stint, αλλαγή ρύθμισης) για να την αποτρέψουν.

\subsubsection{Συνδυασμένο μοντέλο φθοράς}

Και οι τρεις μηχανισμοί δρουν ταυτόχρονα, με τη σχετική τους βαρύτητα να ορίζεται από τη
θερμοκρασία. Για να αποφευχθεί αριθμητική ασυνέχεια στον integrator, η μετάβαση μεταξύ των
μηχανισμών γίνεται μέσω sigmoid γύρω από την κρίσιμη θερμοκρασία $T_{\text{tp}}$:

\begin{equation}
    W_{\text{total}}(T) = W_p + \sigma(T_{\text{tp}} - T) \cdot W_g + \sigma(T - T_{\text{tp}}) \cdot W_b
    \label{eq:Wtotal}
\end{equation}

Η επιλογή sigmoid αντί για step function είναι μεθοδολογικά σημαντική: εξασφαλίζει $C^1$
συνέχεια που επιτρέπει ομαλή αριθμητική ολοκλήρωση. Στην πράξη, η μετάβαση πλάτους $\pm
5\text{--}10\,^\circ\mathrm{C}$ γύρω από $T_{\text{tp}}$ είναι αρκετά απότομη ώστε να
αναπαριστά πιστά το φυσικό φαινόμενο.
